\documentclass[aspectratio=169,xcolor=dvipsnames]{beamer}
\usetheme{Berlin}

\usepackage[spanish]{babel}
\usepackage{hyperref}
\usepackage{graphicx}
\usepackage{booktabs}
\usepackage{amsmath}
\usepackage{amssymb}
\usepackage{lettrine}
\setbeamertemplate{caption}[numbered]
\usepackage[dvipsnames,svgnames,x11names]{xcolor}
\usepackage{xurl}
\usepackage{algorithm}
\usepackage{algorithmicx}
\usepackage{algpseudocode}
\usepackage{adjustbox}
\usepackage{tikz}
\usetikzlibrary{arrows.meta, shapes.geometric, positioning, calc}

\hypersetup{
    colorlinks=true,
    linkcolor=cyan,
    filecolor=blue,
    urlcolor=blue,
    citecolor=blue,
}

%----------------------------------------------------------------------------------------
%   CODE LISTINGS SETTINGS
%----------------------------------------------------------------------------------------
\usepackage{listings}
\definecolor{backcolour}{rgb}{0.97,0.97,0.99}
\definecolor{codegreen}{rgb}{0,0.6,0}

\lstdefinestyle{MATLABStyle}{
  language=Matlab,
  basicstyle=\ttfamily\scriptsize,
  keywordstyle=\color{blue}\bfseries,
  commentstyle=\color{codegreen},
  stringstyle=\color{violet},
  breaklines=true,
  numbers=left,
  numbersep=5pt,
  frame=lines,
  backgroundcolor=\color{backcolour}
}
\lstset{style=MATLABStyle}

%----------------------------------------------------------------------------------------
%   TIKZ STYLES FOR BLOCK DIAGRAMS
%----------------------------------------------------------------------------------------
\tikzset{
  block/.style = {draw, rectangle, minimum height=1.2em, minimum width=2.5em,
                  fill=blue!10, thick},
  sum/.style   = {draw, circle, minimum size=1.2em, thick, fill=yellow!20},
  arrow/.style = {-{Stealth[length=2mm]}, thick},
}

%----------------------------------------------------------------------------------------
%   TITLE CONFIGURATION
%----------------------------------------------------------------------------------------
\title{\'Algebra de Bloques y Representaci\'on en Espacio de Estado}
\subtitle{Simulaci\'on de Sistemas Din\'amicos}
\author{Prof. D.Sc. BARSEKH-ONJI Aboud}
\institute{
    Facultad de Ingenier\'ia \\
    Universidad An\'ahuac M\'exico
}
\date{\today}

%----------------------------------------------------------------------------------------
%   AUTO AGENDA AT EACH SECTION
%----------------------------------------------------------------------------------------
\AtBeginSection[]{
  \begin{frame}{Agenda}
    \tableofcontents[currentsection]
  \end{frame}
}

\begin{document}

\begin{frame}
    \titlepage
\end{frame}

\begin{frame}{Agenda General}
    \tableofcontents
\end{frame}

%================================================
\section{Repaso: Funci\'on de Transferencia}
%================================================

\begin{frame}{?`Qu\'e hemos visto hasta ahora?}
    \begin{block}{Contexto de la clase}
        En sesiones anteriores simulamos sistemas din\'amicos usando tres enfoques en MATLAB:
    \end{block}
    \vspace{0.4em}
    \begin{columns}
        \column{0.33\textwidth}
        \centering
        \textbf{Enfoque 1}\\[0.4em]
        Ecuaciones diferenciales\\
        \texttt{ode45}, \texttt{ode23}
        \column{0.33\textwidth}
        \centering
        \textbf{Enfoque 2}\\[0.4em]
        Diagramas de bloques\\
        Simulink
        \column{0.33\textwidth}
        \centering
        \textbf{Enfoque 3}\\[0.4em]
        \textit{Transfer Function}\\
        \texttt{tf()}, \texttt{step()}, \texttt{lsim()}
    \end{columns}
    \vspace{0.5em}
    \begin{alertblock}{Lo que falta hoy}
        Aprender a \textbf{manipular} diagramas de bloques mediante \'algebra de bloques, y representar sistemas en \textbf{Espacio de Estado}.
    \end{alertblock}
\end{frame}

\begin{frame}{La Funci\'on de Transferencia --- Recordatorio}
    Para un sistema LTI con entrada $U(s)$ y salida $Y(s)$:
    \begin{equation}
        G(s) = \frac{Y(s)}{U(s)} = \frac{b_m s^m + \cdots + b_1 s + b_0}{a_n s^n + \cdots + a_1 s + a_0}
    \end{equation}
    \begin{itemize}
        \item \textbf{Condiciones iniciales:} se asumen cero (Transformada de Laplace unilateral)
        \item \textbf{LTI:} Sistema Lineal e Invariante en el Tiempo
        \item \textbf{Orden del sistema:} determinado por el grado $n$ del denominador
        \item \textbf{Polos:} ra\'ices del denominador $\Rightarrow$ definen la estabilidad
        \item \textbf{Ceros:} ra\'ices del numerador $\Rightarrow$ afectan la respuesta transitoria
    \end{itemize}
\end{frame}

\begin{frame}{Bloque Fundamental}
    Un bloque en un diagrama de control representa una \textbf{funci\'on de transferencia}:
    \vspace{0.6em}
    \begin{center}
    \begin{tikzpicture}
        \node[block, minimum width=3cm, minimum height=1.2cm] (G) {$G(s)$};
        \draw[arrow] (G.west) ++(-1.5,0) node[left]{$U(s)$} -- (G.west);
        \draw[arrow] (G.east) -- ++(1.5,0) node[right]{$Y(s)$};
        \node[below=0.2cm of G] {\footnotesize \textit{Sistema / Planta / Controlador}};
    \end{tikzpicture}
    \end{center}
    \vspace{0.4em}
    \begin{exampleblock}{Relaci\'on algebraica}
        $Y(s) = G(s) \cdot U(s)$ \quad $\Longleftrightarrow$ \quad $G(s) = \dfrac{Y(s)}{U(s)}$
    \end{exampleblock}
\end{frame}

%================================================
\section{\'Algebra de Bloques: Reglas Fundamentales}
%================================================

\begin{frame}{Regla 1 --- Conexi\'on en Serie (Cascada)}
    Cuando la salida de un bloque es la entrada del siguiente:
    \vspace{0.4em}
    \begin{center}
    \begin{tikzpicture}
        \node[block] (G1) {$G_1(s)$};
        \node[block, right=1.8cm of G1] (G2) {$G_2(s)$};
        \draw[arrow] (G1.west) ++(-1.2,0) node[left]{$U(s)$} -- (G1.west);
        \draw[arrow] (G1.east) -- node[above]{\footnotesize $X(s)$} (G2.west);
        \draw[arrow] (G2.east) -- ++(1.2,0) node[right]{$Y(s)$};
    \end{tikzpicture}
    \end{center}
    \vspace{0.3em}
    \begin{equation}
        G_{eq}(s) = G_1(s) \cdot G_2(s)
    \end{equation}
    \begin{block}{Demostraci\'on}
        $X(s) = G_1(s)\,U(s)$ \quad y \quad $Y(s) = G_2(s)\,X(s)$\\[0.3em]
        $\therefore \quad Y(s) = G_2(s)\cdot G_1(s)\,U(s)$
    \end{block}

\end{frame}
\begin{frame}{Regla 1 --- Conexi\'on en Serie (Cascada)}

    \begin{alertblock}{Cuidado}
        Para sistemas SISO: $G_1 G_2 = G_2 G_1$. En sistemas MIMO (matriciales) esto no es v\'alido en general.
    \end{alertblock}
\end{frame}

\begin{frame}{Regla 2 --- Conexi\'on en Paralelo}
    Cuando la misma entrada alimenta dos bloques y sus salidas se suman:
    \vspace{0.3em}
    \begin{center}
    \begin{tikzpicture}[node distance=1cm]
        \node[sum] (S) {};
        \node[block, above left=0.9cm and 1.8cm of S] (G1) {$G_1(s)$};
        \node[block, below left=0.9cm and 1.8cm of S] (G2) {$G_2(s)$};
        \coordinate[left=3cm of S] (entry);
        \draw[arrow] (entry) node[left]{$U(s)$} -- ++(0.9,0) coordinate(split);
        \draw[arrow] (split) |- (G1.west);
        \draw[arrow] (split) |- (G2.west);
        \draw[arrow] (G1.east) -| (S.north) node[pos=0.85,right]{\footnotesize$+$};
        \draw[arrow] (G2.east) -| (S.south) node[pos=0.85,right]{\footnotesize$+$};
        \draw[arrow] (S.east) -- ++(1.2,0) node[right]{$Y(s)$};
    \end{tikzpicture}
    \end{center}
    \vspace{0.2em}
    \begin{equation}
        G_{eq}(s) = G_1(s) + G_2(s)
    \end{equation}
\end{frame}

\begin{frame}{Regla 3 --- Retroalimentaci\'on (Lazo Cerrado)}
    La configuraci\'on m\'as importante en control autom\'atico:
    \vspace{0.3em}
    \begin{center}
    \begin{tikzpicture}[node distance=0.8cm]
        \node[sum] (S) {};
        \node[block, right=1.2cm of S] (G) {$G(s)$};
        \node[block, below=0.9cm of G] (H) {$H(s)$};
        \draw[arrow] (S.west) ++(-1.2,0) node[left]{$R(s)$} -- node[above]{\footnotesize$+$} (S.west);
        \draw[arrow] (S.east) -- node[above]{\footnotesize$E(s)$} (G.west);
        \draw[arrow] (G.east) -- ++(1.8,0) node[right]{$Y(s)$} coordinate(out);
        \draw[arrow] (out)+(-0.4,0) |- (H.east);
        \draw[arrow] (H.west) -| (S.south) node[pos=0.9,right]{\footnotesize$-$};
    \end{tikzpicture}
    \end{center}
    \vspace{0.2em}
    \begin{equation}
        \boxed{T(s) = \frac{G(s)}{1 + G(s)\,H(s)}}
    \end{equation}
    \begin{itemize}
        \item \textbf{$G(s)$}: lazo directo (\textit{forward path})
        \item \textbf{$H(s)$}: retroalimentaci\'on; si $H(s)=1$ (unitaria): $T(s) = \dfrac{G(s)}{1+G(s)}$
        \item \textbf{Lazo abierto}: $L(s) = G(s)\,H(s)$
        \item Signo $+$ en retroalimentaci\'on positiva: $T(s) = \dfrac{G(s)}{1 - G(s)\,H(s)}$
    \end{itemize}
\end{frame}

\begin{frame}{Regla 4 --- Nodo de Suma y Punto de Ramificaci\'on}
    \begin{columns}
        \column{0.48\textwidth}
        \begin{block}{Nodo de suma (Comparador)}
            \begin{center}
            \begin{tikzpicture}
                \node[sum] (S) {};
                \draw[arrow] (S.west) ++(-1,0) node[left]{$A$} -- node[above]{\footnotesize$+$} (S.west);
                \draw[arrow] (S.south) ++(0,-0.7) node[below]{$B$} -- node[right]{\footnotesize$-$} (S.south);
                \draw[arrow] (S.east) -- ++(1,0) node[right]{$C=A-B$};
            \end{tikzpicture}
            \end{center}
            Suma algebraica de se\~nales. Los signos se indican expl\'icitamente sobre cada flecha.
        \end{block}
        \column{0.48\textwidth}
        \begin{block}{Punto de ramificaci\'on (\textit{Branch point})}
            \begin{center}
            \begin{tikzpicture}
                \coordinate (bp) at (0,0);
                \fill (bp) circle (2.5pt);
                \draw[arrow] (bp) ++(-1.2,0) node[left]{$X(s)$} -- (bp);
                \draw[arrow] (bp) -- ++(1.2,0) node[right]{$X(s)$};
                \draw[arrow] (bp) -- ++(0,-1.0) node[below]{$X(s)$};
            \end{tikzpicture}
            \end{center}
            La misma se\~nal se distribuye sin p\'erdida a m\'ultiples destinos.
        \end{block}
    \end{columns}
\end{frame}

\begin{frame}{Reglas de Movimiento de Bloques}
    A veces es necesario mover bloques para simplificar el diagrama:
    \vspace{0.3em}
    \begin{table}
        \centering
        \begin{tabular}{lll}
            \toprule
            \textbf{Movimiento} & \textbf{Original} & \textbf{Equivalente} \\
            \midrule
            Mover $G$ antes de nodo suma & $G$ antes, suma despu\'es & Agregar $1/G$ en la rama \\
            Mover $G$ despu\'es de nodo suma & Suma antes, $G$ despu\'es & Agregar $G$ en la rama \\
            Mover $G$ antes de ramificaci\'on & $G$ antes, rama despu\'es & Agregar $1/G$ en la nueva rama \\
            Mover $G$ despu\'es de ramificaci\'on & Rama antes, $G$ despu\'es & Agregar $G$ en la nueva rama \\
            \bottomrule
        \end{tabular}
    \end{table}
\end{frame}

\begin{frame}{Reglas de Movimiento de Bloques}  
    \begin{alertblock}{Principio clave}
        Al mover un bloque, siempre se deben \textbf{compensar} las modificaciones para mantener la equivalencia matem\'atica. El objetivo es transformar el diagrama a formas reconocibles (serie, paralelo, lazo cerrado).
    \end{alertblock}
\end{frame}

%================================================
\section{Ejemplos de Reducci\'on de Diagramas}
%================================================

\begin{frame}{Ejemplo 1 --- Sistema de Primer Orden en Lazo Cerrado}
    \begin{exampleblock}{Problema}
        Hallar $T(s) = Y(s)/R(s)$: \quad $G(s) = \dfrac{10}{s+2}$, \quad $H(s) = 1$ (retroalimentaci\'on unitaria)
    \end{exampleblock}
    \vspace{0.4em}
    \textbf{Soluci\'on:} Aplicando la f\'ormula de lazo cerrado:
    \begin{align}
        T(s) &= \frac{G(s)}{1 + G(s) H(s)} = \frac{\dfrac{10}{s+2}}{1 + \dfrac{10}{s+2}} \notag \\[0.6em]
             &= \frac{\dfrac{10}{s+2}}{\dfrac{s+12}{s+2}} = \frac{10}{s+12}
    \end{align}
    \begin{block}{Interpretaci\'on f\'isica}
        El polo en lazo abierto estaba en $s=-2$. Con retroalimentaci\'on unitaria de ganancia 10, el polo se movi\'o a $s=-12$. El sistema es \textbf{6 veces m\'as r\'apido} y la constante de tiempo pasa de $\tau=0.5\,$s a $\tau=0.083\,$s.
    \end{block}
\end{frame}

\begin{frame}{Ejemplo 2 --- Bloques en Serie con Retroalimentaci\'on}
    \begin{exampleblock}{Datos}
        $G_1(s) = \dfrac{1}{s}$, \quad $G_2(s) = \dfrac{5}{s+1}$, \quad $H(s) = 2$
    \end{exampleblock}
    \vspace{0.2em}
    \textbf{Paso 1:} Reducir la cascada (lazo directo):
    \begin{equation}
        G(s) = G_1(s)\cdot G_2(s) = \frac{1}{s} \cdot \frac{5}{s+1} = \frac{5}{s(s+1)}
    \end{equation}
    \textbf{Paso 2:} Aplicar f\'ormula de retroalimentaci\'on:
    \begin{equation}
        T(s) = \frac{G(s)}{1+G(s)H(s)} = \frac{\dfrac{5}{s(s+1)}}{1 + \dfrac{10}{s(s+1)}} = \frac{5}{s^2+s+10}
    \end{equation}
    \begin{alertblock}{An\'alisis din\'amico}
        $s^2+s+10 = 0 \Rightarrow \omega_n = \sqrt{10} \approx 3.16\,\text{rad/s}$, \quad $\zeta = \dfrac{1}{2\omega_n} \approx 0.158$ (subamortiguado).
    \end{alertblock}
\end{frame}

\begin{frame}{Ejemplo 3 --- Sistema con Dos Lazos Anidados}
    \begin{exampleblock}{Datos}
        $G_c = 10$, $G_1 = \dfrac{1}{s}$, $G_2 = \dfrac{1}{s+5}$, $H_1 = 1$, $H_2 = 1$
    \end{exampleblock}
    \vspace{0.2em}
    \begin{center}
    \begin{tikzpicture}[scale=0.85, transform shape]
        \node[sum] (S1) {};
        \node[block, right=0.7cm of S1] (Gc) {$G_c$};
        \node[sum, right=0.7cm of Gc] (S2) {};
        \node[block, right=0.7cm of S2] (G1) {$G_1$};
        \node[block, right=0.7cm of G1] (G2) {$G_2$};
        \node[block, below=0.5cm of G1] (H1) {$H_1$};
        \node[block, below=0.9cm of Gc] (H2) {$H_2$};
        \draw[arrow] (S1.west) ++(-0.7,0) node[left]{$R$} -- node[above]{\tiny$+$} (S1);
        \draw[arrow] (S1) -- (Gc);
        \draw[arrow] (Gc) -- node[above]{\tiny$+$} (S2);
        \draw[arrow] (S2) -- (G1);
        \draw[arrow] (G1) -- (G2);
        \draw[arrow] (G2.east) -- ++(0.8,0) node[right]{$Y$} coordinate(out);
        \draw[arrow] (out)+(-0.2,0) |- (H1.east);
        \draw[arrow] (H1.west) -| (S2.south) node[pos=0.9,right]{\tiny$-$};
        \draw[arrow] (out)+(-0.2,0) -- ++(0,-1.5) -| (H2.east);
        \draw[arrow] (H2.west) -| (S1.south) node[pos=0.9,right]{\tiny$-$};
    \end{tikzpicture}
    \end{center}
\end{frame}
\begin{frame}{Ejemplo 3 --- Sistema con Dos Lazos Anidados}
    \textbf{Paso 1:} Lazo interno ($G_1$, $H_1=1$):
    \begin{equation}
        G_{int}(s) = \frac{G_1}{1+G_1 H_1} = \frac{1/s}{1+1/s} = \frac{1}{s+1}
    \end{equation}
    \textbf{Paso 2:} Lazo externo --- $G_{ext} = G_c \cdot G_{int} \cdot G_2 = \dfrac{10}{(s+1)(s+5)}$:
    \begin{equation}
        T(s) = \frac{G_{ext}}{1+G_{ext}\cdot H_2} = \frac{10}{s^2+6s+15}
    \end{equation}
\end{frame}

\begin{frame}{Ejemplo 4 --- Paralelo con Retroalimentaci\'on}
    \begin{exampleblock}{Datos}
        $G_1(s) = \dfrac{2}{s+1}$, \quad $G_2(s) = \dfrac{3}{s+3}$, \quad $H(s) = 0.5$
    \end{exampleblock}
    \vspace{0.2em}
    \textbf{Paso 1:} Reducci\'on del paralelo:
    \begin{equation}
        G(s) = G_1(s) + G_2(s) = \frac{2(s+3)+3(s+1)}{(s+1)(s+3)} = \frac{5s+9}{(s+1)(s+3)}
    \end{equation}
    \textbf{Paso 2:} Lazo cerrado:
    \begin{equation}
        T(s) = \frac{G(s)}{1+G(s)H(s)} = \frac{5s+9}{(s+1)(s+3)+0.5(5s+9)} = \frac{5s+9}{s^2+6.5s+7.5}
    \end{equation}
\end{frame}
\begin{frame}{Ejemplo 4 --- Paralelo con Retroalimentaci\'on}
    \begin{block}{Verificaci\'on}
        El grado del denominador (2) es igual al orden del sistema original. \checkmark\\
        Para validar: \texttt{T = feedback(parallel(G1,G2), 0.5)} en MATLAB.
    \end{block}
\end{frame}

\begin{frame}[fragile]{Verificaci\'on en MATLAB --- \'Algebra de Bloques}
    \begin{columns}
        \column{0.55\textwidth}
\begin{lstlisting}[language=Matlab, style=MATLABStyle]
%% Ejemplo 2: Serie + Retroalimentacion
G1 = tf(1, [1 0]);    % 1/s
G2 = tf(5, [1 1]);    % 5/(s+1)
H  = tf(2, 1);        % ganancia H=2

% Lazo directo (cascada)
G  = series(G1, G2);

% Lazo cerrado
T  = feedback(G, H);

disp('Funcion de transferencia:')
T

% Respuesta al escalon
figure;
step(T);
title('Ejemplo 2 - Respuesta al Escalon');
grid on;

%% Ejemplo 4: Paralelo
G1p = tf(2, [1 1]);
G2p = tf(3, [1 3]);
Gp  = parallel(G1p, G2p);
Tp  = feedback(Gp, 0.5);
figure; step(Tp); grid on;
title('Ejemplo 4 - Paralelo');
\end{lstlisting}
        \column{0.42\textwidth}
        \begin{block}{Funciones de MATLAB}
            \begin{itemize}
                \item \texttt{series(G1,G2)}: cascada $G_1 G_2$
                \item \texttt{parallel(G1,G2)}: paralelo $G_1+G_2$
                \item \texttt{feedback(G,H)}: lazo cerrado $-H$
                \item \texttt{feedback(G,H,+1)}: lazo cerrado $+H$
                \item \texttt{step()}: respuesta al escal\'on
                \item \texttt{pzmap()}: polos y ceros
                \item \texttt{pole()}: obtener polos
            \end{itemize}
        \end{block}
    \end{columns}
\end{frame}

%================================================
\section{Introducci\'on al Espacio de Estado}
%================================================

\begin{frame}{Limitaciones de la Funci\'on de Transferencia}
    \begin{columns}
        \column{0.48\textwidth}
        \begin{alertblock}{Desventajas de $G(s)$}
            \begin{itemize}
                \item Solo aplica a sistemas \textbf{LTI}
                \item Dif\'icil para sistemas \textbf{MIMO}
                \item No captura \textbf{estados internos}
                \item Condiciones iniciales limitadas a cero
                \item No permite dise\~no de control \'optimo directo
            \end{itemize}
        \end{alertblock}
        \column{0.48\textwidth}
        \begin{block}{Ventajas del Espacio de Estado}
            \begin{itemize}
                \item Maneja naturalmente sistemas \textbf{MIMO}
                \item Permite analizar \textbf{controlabilidad} y \textbf{observabilidad}
                \item Representa la din\'amica \textbf{interna} completa
                \item Permite condiciones iniciales no cero
                \item Base del control moderno (LQR, filtro de Kalman)
            \end{itemize}
        \end{block}
    \end{columns}
\end{frame}

\begin{frame}{Variables de Estado --- Concepto}
    \begin{block}{Definici\'on}
        Las \textbf{variables de estado} son el conjunto m\'inimo de variables $x_1(t), x_2(t), \ldots, x_n(t)$ que, junto con la entrada $u(t)$ para $t \geq t_0$, determinan completamente el comportamiento futuro del sistema.
    \end{block}
 
    \begin{columns}
        \column{0.48\textwidth}
        \textbf{Ejemplos de variables de estado:}
        \begin{itemize}
            \item \textbf{Mec\'anico:} posici\'on y velocidad $(q, \dot{q})$
            \item \textbf{El\'ectrico:} voltaje en $C$, corriente en $L$
            \item \textbf{T\'ermico:} temperaturas de los nodos
            \item \textbf{Matem\'atico:} la variable y sus derivadas
        \end{itemize}
        \column{0.48\textwidth}
        \begin{exampleblock}{Idea clave}
            \[
                \underbrace{\mathbf{x}(t_0)}_{\text{estado inicial}} + \underbrace{u(t)}_{\text{entrada}} \Rightarrow \underbrace{y(t)}_{\text{salida}}
            \]
            El estado en $t_0$ m\'as la entrada futura determinan completamente la salida futura.
        \end{exampleblock}
    \end{columns}
\end{frame}

\begin{frame}{Representaci\'on General en Espacio de Estado}
    La forma est\'andar de un sistema en espacio de estado es:
    \begin{equation}
        \dot{\mathbf{x}}(t) = A\,\mathbf{x}(t) + B\,\mathbf{u}(t) \qquad \text{(ecuaci\'on de estado)}
    \end{equation}
    \begin{equation}
        \mathbf{y}(t) = C\,\mathbf{x}(t) + D\,\mathbf{u}(t) \qquad \text{(ecuaci\'on de salida)}
    \end{equation}
    \vspace{0.3em}
    \begin{columns}
        \column{0.5\textwidth}
        \begin{itemize}
            \item $\mathbf{x}(t) \in \mathbb{R}^n$: \textbf{vector de estado}
            \item $\mathbf{u}(t) \in \mathbb{R}^m$: \textbf{vector de entrada}
            \item $\mathbf{y}(t) \in \mathbb{R}^p$: \textbf{vector de salida}
        \end{itemize}
        \column{0.5\textwidth}
        \begin{itemize}
            \item $A \in \mathbb{R}^{n\times n}$: \textbf{matriz de sistema}
            \item $B \in \mathbb{R}^{n\times m}$: \textbf{matriz de entrada}
            \item $C \in \mathbb{R}^{p\times n}$: \textbf{matriz de salida}
            \item $D \in \mathbb{R}^{p\times m}$: \textbf{transmisi\'on directa}
        \end{itemize}
    \end{columns}
\end{frame}
\begin{frame}{Representaci\'on General en Espacio de Estado}    
    \begin{alertblock}{Nota}
        Para la mayor\'ia de sistemas f\'isicos $D = 0$ (no hay transmisi\'on directa de entrada a salida).
    \end{alertblock}
\end{frame}

\begin{frame}{De Ecuaci\'on Diferencial a Espacio de Estado}
    \begin{exampleblock}{Procedimiento sistem\'atico}
        Dada: $\quad y^{(n)} + a_{n-1}y^{(n-1)} + \cdots + a_1\dot{y} + a_0 y = b_0 u$
    \end{exampleblock}
    \vspace{0.3em}
    \textbf{Paso 1 --- Definir variables de estado:}
    \begin{equation}
        x_1 = y, \quad x_2 = \dot{y}, \quad x_3 = \ddot{y}, \quad \ldots, \quad x_n = y^{(n-1)}
    \end{equation}
    \textbf{Paso 2 --- Ecuaciones de estado:}
    \[
        \dot{x}_1 = x_2, \quad \dot{x}_2 = x_3, \quad \ldots, \quad \dot{x}_n = -a_0 x_1 - a_1 x_2 - \cdots - a_{n-1} x_n + b_0 u
    \]
    \textbf{Paso 3 --- Ecuaci\'on de salida:} $\quad y = x_1$
\end{frame}
\begin{frame}{De Ecuaci\'on Diferencial a Espacio de Estado}
    \begin{block}{Resultado: Forma Can\'onica Controlable}
        La \'ultima ecuaci\'on $\dot{x}_n$ concentra todos los coeficientes del denominador de $G(s)$. Esta forma se llama \textbf{can\'onica controlable} o \textit{forma compa\~nera}.
    \end{block}
\end{frame}

\begin{frame}{Forma Can\'onica Controlable}
    Para una EDO de orden $n$, la forma can\'onica controlable es:
    \begin{equation}
        A = \begin{bmatrix}
            0 & 1 & 0 & \cdots & 0 \\
            0 & 0 & 1 & \cdots & 0 \\
            \vdots & & & \ddots & \vdots \\
            0 & 0 & 0 & \cdots & 1 \\
            -a_0 & -a_1 & -a_2 & \cdots & -a_{n-1}
        \end{bmatrix}, \quad
        B = \begin{bmatrix} 0 \\ 0 \\ \vdots \\ 0 \\ 1 \end{bmatrix}
    \end{equation}
    \begin{equation}
        C = \begin{bmatrix} b_0 & b_1 & \cdots & b_{n-1} \end{bmatrix}, \quad D = [0]
    \end{equation}
\end{frame}
\begin{frame}{Forma Can\'onica Controlable}

    \begin{exampleblock}{Propiedad importante}
        Los \textbf{eigenvalores} de la matriz $A$ son los \textbf{polos} del sistema, es decir, las ra\'ices del polinomio caracter\'istico $\det(sI - A)$.
    \end{exampleblock}
\end{frame}

%================================================
\section{Ejemplos en Espacio de Estado}
%================================================

\begin{frame}{Ejemplo 5 --- Sistema de Segundo Orden}
    \begin{exampleblock}{Problema}
        Convertir a espacio de estado: $\quad \ddot{y} + 3\dot{y} + 2y = u$
    \end{exampleblock}
    \vspace{0.3em}
    \textbf{Variables de estado:} $x_1 = y$, $x_2 = \dot{y}$
    \vspace{0.3em}
    \textbf{Forma matricial:}
    \begin{equation}
        \underbrace{\begin{bmatrix} \dot{x}_1 \\ \dot{x}_2 \end{bmatrix}}_{\dot{\mathbf{x}}}
        =
        \underbrace{\begin{bmatrix} 0 & 1 \\ -2 & -3 \end{bmatrix}}_{A}
        \underbrace{\begin{bmatrix} x_1 \\ x_2 \end{bmatrix}}_{\mathbf{x}}
        +
        \underbrace{\begin{bmatrix} 0 \\ 1 \end{bmatrix}}_{B} u, \qquad
        y = \underbrace{\begin{bmatrix} 1 & 0 \end{bmatrix}}_{C} \mathbf{x}
    \end{equation}
\end{frame}
\begin{frame}{Ejemplo 5 --- Sistema de Segundo Orden}
    \begin{block}{Verificaci\'on: $G(s)$ desde Espacio de Estado}
        \[
            G(s) = C(sI-A)^{-1}B = \begin{bmatrix}1&0\end{bmatrix}
            \begin{bmatrix}s+3 & 1 \\ -2 & s\end{bmatrix}^{-1} \frac{1}{s^2+3s+2}
            \begin{bmatrix}0\\1\end{bmatrix} = \frac{1}{s^2+3s+2} \checkmark
        \]
    \end{block}
\end{frame}

\begin{frame}{Ejemplo 6 --- Sistema de Tercer Orden}
    \begin{exampleblock}{Problema}
        $\dddot{y} + 6\ddot{y} + 11\dot{y} + 6y = 2u$
    \end{exampleblock}
    \vspace{0.2em}
    \textbf{Variables:} $x_1 = y$, $x_2 = \dot{y}$, $x_3 = \ddot{y}$
    \begin{equation}
        A = \begin{bmatrix} 0 & 1 & 0 \\ 0 & 0 & 1 \\ -6 & -11 & -6 \end{bmatrix}, \quad
        B = \begin{bmatrix} 0 \\ 0 \\ 2 \end{bmatrix}, \quad
        C = \begin{bmatrix} 1 & 0 & 0 \end{bmatrix}, \quad D = [0]
    \end{equation}
\end{frame}
\begin{frame}{Ejemplo 6 --- Sistema de Tercer Orden}
    \begin{block}{Eigenvalores de $A$ (polos del sistema)}
        \[
            \det(\lambda I - A) = \lambda^3 + 6\lambda^2 + 11\lambda + 6 = (\lambda+1)(\lambda+2)(\lambda+3) = 0
        \]
        $\lambda_1 = -1$, \; $\lambda_2 = -2$, \; $\lambda_3 = -3$ \quad $\Rightarrow$ \textbf{sistema estable} (todos en semiplano izquierdo).
    \end{block}
\end{frame}

\begin{frame}{Ejemplo 7 --- Circuito RLC en Espacio de Estado}
    \begin{columns}
        \column{0.45\textwidth}
        \begin{exampleblock}{Sistema el\'ectrico}
            Circuito RLC serie con $R=1\,\Omega$, $L=1\,H$, $C=0.5\,F$.\\[0.5em]
            Ecuaci\'on de malla:
            \[
                L\ddot{i} + R\dot{i} + \frac{i}{C} = \dot{u}
            \]
            En t\'erminos de $v_C$:
            \[
                \ddot{v}_C + \dot{v}_C + 2v_C = 2u
            \]
        \end{exampleblock}
        \column{0.52\textwidth}
        \textbf{Variables de estado f\'isicas:}
        \begin{itemize}
            \item $x_1 = v_C$ (voltaje en capacitor)
            \item $x_2 = i_L = C\dot{v}_C$ (corriente en inductor)
        \end{itemize}
        \vspace{0.3em}
        \begin{equation}
            A = \begin{bmatrix} 0 & 2 \\ -0.5 & -1 \end{bmatrix}, \quad
            B = \begin{bmatrix} 0 \\ 1 \end{bmatrix}
        \end{equation}
        \begin{equation}
            C = \begin{bmatrix} 1 & 0 \end{bmatrix}, \quad D = [0]
        \end{equation}
        \vspace{0.2em}
        \textit{Las variables de estado tienen \textbf{significado f\'isico} directo.}
    \end{columns}
\end{frame}

\begin{frame}[fragile]{Espacio de Estado en MATLAB}
    \begin{columns}
        \column{0.55\textwidth}
\begin{lstlisting}[language=Matlab, style=MATLABStyle]
%% Ejemplo 5 - Sistema 2do orden
A = [0  1; -2 -3];
B = [0; 1];
C = [1  0];
D = 0;

% Crear sistema en espacio de estado
sys_ss = ss(A, B, C, D);

% Convertir a funcion de transferencia
sys_tf = tf(sys_ss);
disp('G(s):'); sys_tf

% Respuesta al escalon
figure; step(sys_ss); grid on;
title('Respuesta al Escalon (EE)');

% Eigenvalores = Polos
disp('Eigenvalores:'); eig(A)
disp('Polos:'); pole(sys_ss)
\end{lstlisting}
        \column{0.42\textwidth}
\begin{lstlisting}[language=Matlab, style=MATLABStyle]
%% Desde TF hacia EE
G = tf(1, [1 3 2]);
sys2 = ss(G);  % conversion
sys2.A
sys2.B
sys2.C
sys2.D

%% Simulacion con CI no cero
x0 = [1; 0];      % condicion inicial
t  = 0:0.01:10;
u  = zeros(size(t));
[y, tout, x] = lsim(sys2, u, t, x0);
figure;
plot(tout, y, 'b', ...
     tout, x(:,1), 'r--');
legend('y(t)', 'x_1(t)');
xlabel('t (s)');
title('Respuesta libre con CI');
\end{lstlisting}
    \end{columns}
\end{frame}

%================================================
\section{S\'intesis y Conexi\'on entre Representaciones}
%================================================

\begin{frame}{Relaci\'on entre EE y Funci\'on de Transferencia}
    La funci\'on de transferencia se obtiene a partir de las matrices del EE:
    \begin{equation}
        \boxed{G(s) = C\,(sI - A)^{-1}\,B + D}
    \end{equation}
    \vspace{0.3em}
    \begin{columns}
        \column{0.48\textwidth}
        \begin{block}{De $G(s)$ a Espacio de Estado}
            \begin{enumerate}
                \item Identificar coeficientes del numerador y denominador
                \item Aplicar forma can\'onica controlable (u observable)
                \item Construir matrices $A, B, C, D$
                \item En MATLAB: \texttt{ss(G)}
            \end{enumerate}
        \end{block}
        \column{0.48\textwidth}
        \begin{block}{De Espacio de Estado a $G(s)$}
            \begin{enumerate}
                \item Calcular $(sI-A)^{-1}$ via adjunta/determinante
                \item Multiplicar $C \cdot (sI-A)^{-1} \cdot B$
                \item Sumar $D$ si hay transmisi\'on directa
                \item En MATLAB: \texttt{tf(sys\_ss)}
            \end{enumerate}
        \end{block}
    \end{columns}
    \vspace{0.3em}
    \begin{alertblock}{Importante: No unicidad}
        La representaci\'on en EE \textbf{no es \'unica}. Existen infinitas representaciones para la misma $G(s)$, pero todas comparten los mismos eigenvalores (polos).
    \end{alertblock}
\end{frame}

\begin{frame}{Tabla Comparativa de Representaciones}
    \begin{table}
        \centering
        \begin{tabular}{lccc}
            \toprule
            \textbf{Caracter\'istica} & \textbf{EDO} & \textbf{$G(s)$} & \textbf{$(A,B,C,D)$} \\
            \midrule
            Sistemas SISO          & \checkmark & \checkmark & \checkmark \\
            Sistemas MIMO          & Complejo   & Matricial  & \checkmark \\
            No lineales            & \checkmark & $\times$   & Parcial    \\
            Condiciones iniciales  & \checkmark & $\times$   & \checkmark \\
            An\'alisis de estabilidad & Moderado & F\'acil   & \checkmark \\
            Controlabilidad/Observ. & $\times$  & Limitado   & \checkmark \\
            Dise\~no de control    & Dif\'icil  & Cl\'asico  & Moderno    \\
            MATLAB                 & \texttt{ode45} & \texttt{tf} & \texttt{ss} \\
            \bottomrule
        \end{tabular}
        \caption{Comparaci\'on de representaciones de sistemas din\'amicos}
    \end{table}
\end{frame}

\begin{frame}{Resumen de Reglas del \'Algebra de Bloques}
    \begin{table}
        \centering
        \begin{tabular}{lll}
            \toprule
            \textbf{Configuraci\'on} & \textbf{Diagrama} & \textbf{Equivalente} \\
            \midrule
            Serie (cascada)       & $G_1 \rightarrow G_2$           & $G_1 \cdot G_2$ \\[0.3em]
            Paralelo              & $G_1 \parallel G_2$             & $G_1 + G_2$ \\[0.3em]
            Lazo cerrado ($-$)    & $G$ con retroalim. $H$          & $\dfrac{G}{1+GH}$ \\[0.5em]
            Lazo cerrado ($+$)    & $G$ con retroalim. $+H$         & $\dfrac{G}{1-GH}$ \\[0.5em]
            Retroalim. unitaria   & $G$, $H=1$                      & $\dfrac{G}{1+G}$ \\
            \bottomrule
        \end{tabular}
        \caption{Reglas fundamentales de reducci\'on de diagramas de bloques}
    \end{table}

\end{frame}
\begin{frame}{Resumen de Reglas del \'Algebra de Bloques}
\begin{alertblock}{Estrategia}
        Reducir primero los \textbf{lazos internos}, luego los externos. Simplificar serie y paralelo antes de aplicar la f\'ormula de retroalimentaci\'on.
    \end{alertblock}
\end{frame}

\begin{frame}{Actividad Pr\'actica}
    \begin{exampleblock}{Ejercicio}
        Dado el sistema: $\quad \ddot{y} + 5\dot{y} + 6y = 3u$
    \end{exampleblock}
    Realizar en MATLAB:
    \begin{enumerate}
        \item Representar como \textbf{Funci\'on de Transferencia} con \texttt{tf()}
        \item Representar en \textbf{Espacio de Estado} con \texttt{ss()} (desde la EDO y desde \texttt{tf})
        \item Verificar que ambas dan la misma respuesta al escal\'on con \texttt{step()}
        \item Calcular los \textbf{polos} con \texttt{pole()} y los \textbf{eigenvalores} con \texttt{eig(A)}
        \item Agregar un \textbf{controlador proporcional} $G_c = K = 5$ en lazo cerrado con \texttt{feedback()}
        \item Comparar la respuesta al escal\'on en lazo \textbf{abierto} vs \textbf{cerrado}
        \item Entregables: Script de MATLAB comentado con figuras de respuesta al escal\'on para los puntos 3 y 6.
    \end{enumerate}

\end{frame}

\begin{frame}{Referencias y Recursos}
    \begin{exampleblock}{Bibliograf\'ia recomendada}
        \begin{itemize}
            \item Ogata, K. (2010). \textit{Modern Control Engineering} (5th ed.). Prentice Hall.
            \item Nise, N. S. (2019). \textit{Control Systems Engineering} (8th ed.). Wiley.
            \item Franklin, G., Powell, J. D., \& Emami-Naeini, A. (2019). \textit{Feedback Control of Dynamic Systems}. Pearson.
        \end{itemize}
    \end{exampleblock}
    \vspace{0.3em}
    \begin{exampleblock}{Documentaci\'on MATLAB}
        \begin{itemize}
            \item \url{https://www.mathworks.com/help/control/}
            \item Funciones: \texttt{tf}, \texttt{ss}, \texttt{series}, \texttt{parallel}, \texttt{feedback}, \texttt{step}, \texttt{lsim}, \texttt{pzmap}, \texttt{eig}, \texttt{pole}
        \end{itemize}
    \end{exampleblock}
    \vspace{0.5em}
    \begin{center}
        \Large \textbf{!`Gracias!}\\[0.3em]
        \normalsize \textit{Prof. D.Sc. BARSEKH-ONJI Aboud}\\
        \small Facultad de Ingenier\'ia --- Universidad An\'ahuac M\'exico
    \end{center}
\end{frame}

\end{document}
